% Corrigé intégré automatiquement pour la banque Exercices.
% Source : A_Integrer/DDS_04/072_SLCI_PI/072_SLCI_PI.tex
\begin{corrigebox}[Corrigé]
\CorrigeQuestion{1}
On réduit le schéma-blocs de l'intérieur vers l'extérieur : associations en série par produit, associations en parallèle par somme et boucle de retour par $H_{BF}=\dfrac{G}{1+GH}$ (retour négatif). Après simplification, on ordonne le numérateur et le dénominateur selon les puissances de $p$, puis on identifie le gain statique, la classe et les paramètres canoniques.
\CorrigeQuestion{2}
On applique le théorème de la valeur finale : $y(\infty)=\lim_{p\to0}pY(p)$, après avoir vérifié la stabilité de $pY(p)$. L'écart permanent se déduit de $E(p)=\dfrac{R(p)}{1+L(p)}$ en tenant compte de la classe de la boucle ouverte et de la nature de l'entrée.
\CorrigeQuestion{3}
$F_2(p)=\dfrac{K_S}{\dfrac{p^2}{\omega_1^2}+pk_v K_S +k_pK_S - 1}$.
\CorrigeQuestion{4}
$k_p>\dfrac{1}{K_S}$ et $k_v>0$.
\CorrigeQuestion{5}
$K_2 = \dfrac{K_S}{k_p K_S - 1}$, $\omega_0 = \omega_1\sqrt{k_p K_S - 1}$ et $\xi = \dfrac{1}{2}\dfrac{k_vK_S\omega_1}{\sqrt{k_p K_S - 1}}$.
\CorrigeQuestion{6}
%%
%%
%%
\end{corrigebox}
